\section{Conclusion}
\label{sec:conclusion}
We presented \textbf{InfoDelphi}, a multi-agent forecasting framework built on a simple but overlooked insight: collective reasoning only outperforms individual reasoning when participants hold genuinely different information.
Treating information asymmetry as a first-class design principle --- and correcting overconfidence in the aggregation layer rather than the prompt --- yields the best Brier score among ten systems on two fixed-evidence benchmarks, with paired-bootstrap significance against every baseline.
Our analysis localizes the mechanism (herding collapse under homogeneous input) and its boundary (evidence-limited domains), suggesting that in multi-agent LLM systems, \emph{what each agent reads} deserves the same design attention as \emph{how agents talk}.

\section*{Limitations}

\paragraph{Theoretical assumptions.}
Proposition~\ref{prop:decorrelation} rests on three idealizing assumptions: a linear error decomposition, public--private independence, and symmetric private noise. In practice, documents retrieved for the same question share topical correlations, partially violating independence and making the decorrelation bound approximate; LLMs, as non-linear functions of their inputs, need not satisfy the fixed-weight decomposition.
Consistently, the empirical error correlations in Figure~\ref{fig:herding} remain high (${\sim}0.94$) even with partitioned evidence, because model-level correlation persists regardless of evidence allocation.

\paragraph{Effect sizes and scope.}
The deliberation-round effect is small relative to question-level noise and not individually significant, and the advantage over the strongest homogeneous-input baseline, while significant on \textsc{PolyGym}, is directional on FutureX; when the shrinkage transform is granted to all baselines (Appendix~\ref{sec:appendix_control}), InfoDelphi remains best on both datasets but several pairwise gaps are no longer individually significant.
Our primary experiments use a single backbone model (\texttt{gpt-5.4-mini}), binary questions, and fixed evidence pools; live-retrieval settings, open-ended questions, and heterogeneous model panels remain untested.
The shrinkage parameters, while validated across our two datasets, are not guaranteed to transfer to models or question distributions with different bias profiles.
API cost scales linearly with agents $\times$ rounds ($1{,}500$ forecasting calls for a full \textsc{PolyGym}-250 run at $J{=}3$, $R{=}2$), which limited exploration of larger panels and longer deliberation.

\section*{Ethics Considerations}
This work presents a forecasting framework and does not directly enable disinformation or harm beyond risks inherent to LLM-based systems generally. However, automated probability estimates generated at scale could in principle be misused to manipulate market sentiment or create false impressions of consensus. We recommend that InfoDelphi be used as a decision-support tool with human oversight in high-stakes deployment contexts.
